\chapter{Proposed Design}
\label{ch:design_problem}

The two preceding chapters established \emph{why} adaptive, safety-bounded load
management is worth building and \emph{where} the existing market and research
leave a gap. This chapter and the next begin the engineering answer. They form a
single Proposed Design chapter split into two halves. The present half states the
problem formally and derives, from that statement, the standards the design must
honour, the wider contexts it must answer to, the requirements and specifications
it must meet, the constraints it operates under, and the risks it sets out to
manage. The second half (\cref{ch:design_solution}) takes those requirements and
develops the concrete architecture, namely the three-plane separation, the
cross-service contracts, the routing safety hierarchy, and the
deterministic-fallback discipline that satisfies them. We keep the present
half requirements-focused and defer the large architecture diagrams to
\cref{ch:design_solution}, so that the solution can be read as a direct response
to a problem that has already been pinned down.

%======================================================================
\section{Problem Definition and Project Purpose}
\label{dsa:sec:problem}

\subsection{Problem statement}
\label{dsa:sec:problem-statement}

A pool of interchangeable backend services behind a single load balancer is the
standard architecture for serving a web application, and it has a standard limit.
The load balancer distributes requests across the pool using a fixed,
deterministic rule, typically round-robin or least-connections, and the size
of the pool is held constant or adjusted only \emph{after} a monitoring threshold
is breached. Both decisions are blind to where demand is heading and to the live
health of the individual backends. When demand rises faster than a reactive
autoscaler can provision new capacity, requests queue and users experience the
peak as latency or errors. When one backend in the pool degrades, whether through
a slow disk, a warm-up stall, or a noisy co-tenant on a shared host, a blind
scheduler keeps sending it its share of traffic, and every request unlucky enough
to land there is penalised. The conventional remedy, permanently over-provisioning
the pool so that peaks and sick instances are absorbed by slack, trades a
reliability problem for a cost-and-energy problem: capacity is paid for and powered
continuously to serve load that arrives only occasionally.

The reactive limb of that limit is worth making concrete, because it is structural
rather than a tuning oversight. A threshold-driven autoscaler can only act once a
metric has already crossed its trigger, after which a stabilisation window guards
against flapping and only then does the replica count move. The new instance is not
useful the moment it is requested either: it must start its container and warm its
caches before it can carry traffic. \Cref{dsa:fig:lag} traces this sequence on a
single time axis. The lag between load arriving and fresh capacity becoming useful
is bounded below by container start plus warm-up, commonly on the order of
30 to 90 seconds for typical web workloads \cite{verma2015borg}. By the
time the new backend is serving, the spike it was provisioned for has already been
absorbed as user-visible latency. Reacting after the alarm is, by construction, too
late; the only way to have capacity ready when the spike lands is to commit to it
\emph{before} the threshold is crossed, which is what a forecast-driven scaler sets
out to do.

\begin{figure}[ht]
\centering
\begin{tikzpicture}[x=1.05cm,y=1cm]
  % time axis
  \draw[slbus,-{Stealth[length=3mm]}] (0,0) -- (12.6,0)
    node[below=5pt,font=\footnotesize\itshape] {time};
  \foreach \x in {0,1.6,4.6,7.6,9.6,11.4} { \draw[slmuted] (\x,0.1) -- (\x,-0.1); }

  % demand curve (the spike) sitting above the axis
  \draw[thick,slbase]
    plot[smooth,tension=0.7] coordinates
    {(0,0.55) (1.6,0.65) (3.0,2.55) (4.6,3.05) (6.0,2.55) (7.6,1.45) (9.6,0.8) (11.4,0.65)};
  \node[font=\footnotesize\itshape,slbase] at (4.6,3.45) {demand};

  % event markers along the axis
  \node[slnode,fill=slbox,font=\scriptsize,text width=1.9cm,align=center,anchor=north]
    at (1.6,-0.45) {load begins to rise};
  \node[slnode,fill=slhead,font=\scriptsize,text width=2.2cm,align=center,anchor=north]
    at (4.6,-0.45) {threshold breached;\\ stabilisation wait};
  \node[slnode,fill=slhead,font=\scriptsize,text width=2.1cm,align=center,anchor=north]
    at (7.6,-0.45) {container start\\ \texttt{+} warm-up};
  \node[slnode,fill=slmint,font=\scriptsize,text width=2.1cm,align=center,anchor=north]
    at (10.5,-0.45) {capacity ready\\ \& serving};

  % the lag span (reaction interval)
  \draw[decorate,decoration={brace,amplitude=5pt,mirror},slmuted]
    (4.6,-2.5) -- (10.5,-2.5)
    node[midway,below=6pt,font=\scriptsize,align=center,text width=5cm]
    {reaction lag bounded below by\\ container start \texttt{+} warm-up
     (typically 30--90\,s)};

  % the "already hurt" region
  \draw[densely dashed,tbdred] (4.6,0) -- (4.6,3.05);
  \draw[densely dashed,tbdred] (10.5,0) -- (10.5,-1.9);
  \node[font=\scriptsize\itshape,tbdred,align=center,text width=3.2cm]
    at (6.6,1.9) {spike served as\\ user-visible latency\\ during the lag};
\end{tikzpicture}
\caption[The structural lag in reactive scaling]{The structural lag in reactive scaling: the interval between a threshold breach and useful capacity is absorbed as degraded service.}
\label{dsa:fig:lag}
\end{figure}

The engineering problem this project addresses is therefore the following: \emph{how
can the routing and capacity decisions at the front door of a backend pool be made
adaptive, driven by live telemetry and short-horizon demand forecasts and able
to learn from backend behaviour, without giving up the deterministic
guarantees that make a classical load balancer safe to put in the request path?}
The tension is built in. The very adaptiveness that lets a system act ahead of a
surge, or steer traffic away from a sick instance, also introduces a learned
component that can be wrong, and a wrong decision in the request path is not a
wrong answer in a report. It is live user traffic sent to the wrong place. A
design that resolves this tension must let the adaptive parts be exercised only
inside bounds that a deterministic baseline and a human operator can always
re-assert.

\subsection{Project purpose}
\label{dsa:sec:problem-purpose}

The purpose of SmartLoad is to deliver that adaptive front door as a single,
self-hosted middleware unit that a small engineering team can run itself. It
unifies, behind one deployable stack, the concerns that are otherwise bought and
operated as separate tools: a load-balancing data plane, anomaly detection over
backend health, short-horizon demand forecasting, proactive and bidirectional
autoscaling, learned request routing, a centralised policy and audit surface for
operators, and the telemetry pipeline that feeds every one of those decisions. The
unifying design commitment, the property that sets SmartLoad apart from a
loose assembly of those tools, is that each adaptive decision sits on top of a
deterministic fallback and under an operator override. The canonical project
definition records this directly: loss of any intelligent component must reduce the
system to classical routing, never to an outage. That commitment is what makes the
learned parts defensible in the request path, and it is the thread this chapter
follows from problem to requirement to constraint.

That commitment resolves into four design principles that the requirements and the
architecture both answer to. \Cref{dsa:fig:principles} states them together because
they are mutually supporting rather than independent. A \emph{deterministic floor}
under every adaptive decision is only safe to rely on if the request path is also
\emph{stateless}, so that any backend can serve any request and the floor can reroute
freely; the floor is only swappable for a learned upgrade because every engine sits
behind a uniform \emph{plugin contract}; and the combination of a floor, a stateless
path, and a contract is what delivers \emph{graceful degradation}, the property that
losing an intelligent component costs adaptation rather than availability. The four
recur throughout the chapter: the deterministic floor as the baseline column of
\cref{dsa:tab:fr}, graceful degradation as NFR-1, the plugin contract as NFR-6, and
the stateless path as the first bounding assumption of \cref{dsa:sec:problem-scope}.

\begin{figure}[ht]
\centering
\begin{tikzpicture}[
  font=\footnotesize,
  prin/.style={slnode,fill=slbox,text width=3.5cm,align=center,
               minimum height=1.9cm,inner sep=5pt,rounded corners=2pt},
  core/.style={slnode,fill=slmint,text width=2.6cm,align=center,
               minimum height=1.5cm,inner sep=4pt},
  link/.style={slmuted,-{Stealth[length=2mm]}},
]
  \node[core] (c) at (0,0) {Adaptive decision\\ under deterministic\\ bound};
  \node[prin] (p1) at (-4.2, 2.0)
    {\textbf{Deterministic floor}\\[2pt] a baseline backs every engine and cannot
     raise};
  \node[prin] (p2) at ( 4.2, 2.0)
    {\textbf{Graceful degradation}\\[2pt] losing an engine stops adaptation, not
     traffic};
  \node[prin] (p3) at (-4.2,-2.0)
    {\textbf{Stateless request path}\\[2pt] any backend serves any request, so
     rerouting is free};
  \node[prin] (p4) at ( 4.2,-2.0)
    {\textbf{Plugin contract}\\[2pt] each engine is swappable behind one typed
     interface};
  \draw[link] (c) -- (p1);
  \draw[link] (c) -- (p2);
  \draw[link] (c) -- (p3);
  \draw[link] (c) -- (p4);
\end{tikzpicture}
\caption[Four design principles for a safe adaptive path]{The four mutually supporting design principles that make an adaptive decision safe to place in the request path.}
\label{dsa:fig:principles}
\end{figure}

\subsection{Scope: what is solved and what is deferred}
\label{dsa:sec:problem-scope}

SmartLoad is scoped, in its present phase, as a \emph{single-tenant, self-hosted
adaptive middleware} for the request path of a homogeneous backend pool. Within
that scope it is responsible, end to end, for the full adaptive control loop:
observing per-request telemetry from the data plane, detecting unhealthy backends,
forecasting demand, deciding how to route and whether to scale, and actuating those
decisions on the load balancer, all under an operator-set policy with a complete
audit trail. \Cref{dsa:tab:scope} fixes the boundary explicitly, separating what
the prototype solves from what it deliberately defers.

\begin{table}[ht]
\centering
\caption[Project scope and Phase 2 deferrals]{Project scope, where the right-hand column lists deliberate Phase~2 deferrals rather than unfinished work.}
\label{dsa:tab:scope}
\small
\begin{tabular}{@{}p{0.46\linewidth}p{0.46\linewidth}@{}}
\toprule
\textbf{In scope (present phase)} & \textbf{Out of scope (deferred to Phase~2)} \\
\midrule
HTTP/1.1 reverse proxying on an NGINX data plane &
Multi-tenant policy and telemetry isolation \\
Classical load balancing (round-robin, least-connections, weighted) &
Operator authentication, role-based access control (RBAC) \\
Telemetry ingest over OTLP, storage in TimescaleDB &
Per-token rate limiting and quotas \\
Anomaly detection on backend health, with backend exclusion &
TLS termination at the load balancer \\
Short-horizon demand forecasting + proactive autoscaling &
Outbound webhook dispatcher for external integrators \\
Learned (RL) routing in shadow then active mode &
A managed, multi-customer SaaS control plane \\
Centralised policy manager with a single safe-mode flag &
Persistent control bus (Redis is treated as in-memory only) \\
Containerised single-node deployment via Docker Compose &
Production-grade Kubernetes operator and HPA integration \\
\bottomrule
\end{tabular}
\end{table}

The deferrals share a theme: every item in the right-hand column belongs to turning
a self-hosted middleware into a managed, multi-customer service. Excluding them in
this phase is what allows the present effort to demonstrate the core adaptive
control loop end to end, the genuinely novel and demanding part, rather than
spreading thin across a product surface whose hard problems are well understood.
The work also rests on four bounding assumptions that the scope makes explicit: the
request path carries no per-client session state (no sticky sessions), so any
backend can serve any request; the control bus may lose messages and is treated as
best-effort transport, so every consumer of an adaptive signal must tolerate its
absence; backends are interchangeable members of a pool rather than functionally
distinct services; and the operator, not the system, is the final authority over
every decision.

%======================================================================
\section{Selected Standards and How They Are Used in the Design}
\label{dsa:sec:standards}

A middleware that sits between a data plane, an observability stack, and a set of
operator tools earns its portability by speaking established standards at every
interface where one exists, rather than inventing private formats that lock
adopters in. The design therefore treats standards as binding constraints on its
external surfaces, not as documentation written after the fact. This section names
the standards SmartLoad adopts and gives concrete evidence of where each is
incorporated; \cref{dsa:tab:standards} summarises the mapping, and the paragraphs
that follow expand on the load-bearing cases.

\begin{table}[ht]
\centering
\caption[Standards adopted and their design artefacts]{Standards and conventions adopted, and the concrete design artefact that incorporates each.}
\label{dsa:tab:standards}
\small
\begin{tabular}{@{}p{0.30\linewidth}p{0.30\linewidth}p{0.32\linewidth}@{}}
\toprule
\textbf{Standard / convention} & \textbf{Where it is used} & \textbf{Evidence of incorporation} \\
\midrule
OpenAPI 3.1 \cite{openapi} & The operator / integrator HTTP control surface &
\texttt{docs/openapi/smartload-v1.yaml} specifies every shipped \texttt{/api/v1/*} route; a structural lint checks code against it \\
OpenTelemetry / OTLP \cite{opentelemetry} & Request-telemetry transport from the data plane &
The load-balancer log shipper emits OTLP/HTTP-JSON to a collector that forwards to the telemetry service \\
Prometheus exposition format \cite{prometheus} & Per-process service health metrics &
Each control-plane service exposes a \texttt{/metrics} endpoint scraped by Prometheus and rendered in Grafana \cite{grafana} \\
RESTful API conventions & Resource modelling of the HTTP surface &
Resource-oriented paths, HTTP verbs, and status-code semantics (200 / 400 / 404 / 503) throughout the spec \\
Semantic versioning & API, policy schema, Redis envelope, client SDK &
A breaking change requires a path-version bump (\texttt{/api/v2}) plus a \texttt{Sunset} header on the deprecated version \\
Redis pub/sub envelope contract & Inter-service control-plane messaging &
Five named channels, each with a typed, versioned envelope, catalogued in \texttt{docs/redis-channels.md} \\
\bottomrule
\end{tabular}
\end{table}

\subsection{OpenAPI 3.1 for the HTTP control surface}
\label{dsa:sec:standards-openapi}

Every operator and integrator interaction with SmartLoad that is not a fire-and-forget
event happens over HTTP, and that surface is specified against OpenAPI~3.1
\cite{openapi} in a single canonical document, \texttt{docs/openapi/smartload-v1.yaml}.
The specification is not aspirational: it enumerates each shipped \texttt{/api/v1/*}
route, namely policy read and write (\texttt{GET}/\texttt{POST /api/v1/policy}), the
append-only audit reads (\texttt{/api/v1/audit/policy}, \texttt{/api/v1/audit/scaling}),
the manual operator overrides (\texttt{/api/v1/scale}, \texttt{/api/v1/isolate}), the
load-balancer state and override endpoints (\texttt{/api/v1/lb/state},
\texttt{/api/v1/lb/weights}, \texttt{/api/v1/lb/algorithm}), the telemetry reads, and
the consolidated \texttt{/api/v1/status} aggregate, together with their request and
response schemas and error shapes. Because the contract is machine-readable, it
serves three concrete design purposes: the Python client SDK and the operator UI
both consume it as their source of truth, a continuous-integration lint validates
that no route drifts from it, and an external integrator can generate a typed client
without reading SmartLoad's source. The specification also encodes the versioning
discipline described below by declaring that every breaking change requires a path
version bump and a \texttt{Sunset} header on the deprecated version.

\subsection{OpenTelemetry and the Prometheus exposition format for observability}
\label{dsa:sec:standards-telemetry}

SmartLoad adopts a vendor-neutral observability stack rather than emitting telemetry
in a private format, because the value of the decision plane is bounded by the
quality and portability of the telemetry beneath it. Two distinct surfaces use two
distinct standards. Per-\emph{request} telemetry, meaning the per-backend latencies,
status codes, and request counts that the decision engines reason over, is
shipped in OpenTelemetry's OTLP format \cite{opentelemetry}: a log shipper tails the
load balancer's structured access log and posts OTLP/HTTP-JSON to a collector, which
forwards it to the telemetry service for storage in TimescaleDB
\cite{timescaledb}. Adopting OTLP means this stream can flow into any conformant
collector rather than a bespoke agent. Per-\emph{process} health, meaning each
service's own liveness, cycle counters, and decision counters, is published separately in
the Prometheus exposition format \cite{prometheus} on a \texttt{/metrics} endpoint,
scraped by Prometheus and visualised in Grafana \cite{grafana}. Separating the two
surfaces, each on its own standard, keeps the high-volume per-request path
independent of the low-volume per-process path while letting any Prometheus-compatible
scraper ingest service metrics with no custom glue.

\subsection{The Redis pub/sub envelope contract}
\label{dsa:sec:standards-redis}

Where HTTP carries synchronous operator interactions, the inter-service control-plane
messaging is asynchronous and event-driven, and it is governed by an internal but
strictly enforced contract: every decision an engine emits travels as a typed,
versioned envelope on one of exactly five named Redis \cite{carlson2013redis}
channels. These channels, namely \texttt{smartload.policy} (policy updates),
\texttt{smartload.anomaly} (backend health verdicts), \texttt{smartload.forecast}
(demand forecasts), \texttt{smartload.routing} (routing recommendations), and
\texttt{smartload.scale} (scaling actions), are catalogued in
\texttt{docs/redis-channels.md}, which records for each channel its single publisher,
its subscribers, its envelope type, and its publish cadence. The contract is evidence
of incorporation in two senses: the envelopes are defined once as typed structures in
a shared module rather than ad-hoc JSON, and a continuous-integration lint greps the
service source for channel names and refuses any channel that is not catalogued. This
five-channel registry, paired with the OpenAPI document, forms the complete external
and internal contract surface of the system; the channels are detailed as a
first-class architectural element in \cref{ch:design_solution}.

\subsection{REST conventions and semantic versioning}
\label{dsa:sec:standards-rest}

Two further conventions shape the surfaces above. The HTTP API follows RESTful
conventions: resource-oriented paths, conventional verb semantics, and a
disciplined use of status codes (a uniform \texttt{200}/\texttt{400}/\texttt{404}/\texttt{503}
vocabulary, with the consolidated status endpoint deliberately always returning
\texttt{200} and rolling per-service failures into its body rather than failing the
whole call). Semantic versioning is applied consistently across all four contract
artefacts, the HTTP API, the policy schema, the Redis envelope, and the client SDK,
so that a consumer can reason about compatibility from a version number alone, and
so that the obligation to bump a version on a breaking change is uniform rather than
per-surface. Together these conventions make the contract surfaces predictable to an
outside integrator, which is the precondition for the portability concern raised in
\cref{dsa:sec:contexts}.

%======================================================================
\section{Major Concerns in Global, Economic, Environmental, and Societal Contexts}
\label{dsa:sec:contexts}

A component that sits in the request path of a production service is never a purely
technical artefact; design choices in it have consequences that reach beyond the
code. This section names the concerns in each of the four contexts that the rubric
requires and, crucially, traces each one to a concrete design pressure rather than
leaving it as a general sentiment. The same concerns are taken up quantitatively
where evidence exists in \cref{ch:discussion}; here we establish how they shape the
requirements that the rest of this chapter derives.

\paragraph{Economic context.}
The dominant economic concern the design answers is the cost of standing capacity.
The conventional way to survive demand spikes is to leave spare backends permanently
running so that slack absorbs the peak; that slack is paid for continuously to serve
load that arrives only occasionally. SmartLoad's economic purpose is to replace
standing slack with \emph{right-sized} compute, growing the pool ahead of a
forecast peak and, equally importantly, shrinking it again when demand recedes. This
concern is the direct origin of the requirement that the autoscaler be bidirectional
and forecast-driven rather than merely reactive (\cref{dsa:sec:reqs-functional}). A
second economic concern shapes the product's position: the design targets the team
below the threshold at which a full orchestration platform is justified, so the
choice of a single self-hosted unit that installs from a container manifest is itself
an economic design decision aimed at a customer who cannot justify standing up and
staffing a Kubernetes control plane.

\paragraph{Environmental context.}
Right-sizing compute is at the same time an environmental concern. Idle servers held
in reserve draw power and cooling to serve a load that is not present; capacity that
is scaled down and switched off is energy that is not wasted. The environmental
dividend is therefore not a separate feature but the same bidirectional-scaling
mechanism viewed through a different lens: the autoscaler's commitment to scaling the
pool \emph{down} as demand subsides, not only up as it rises, is an
energy-footprint decision as much as a cost decision. At the scale of one small
deployment the saving is modest; as a default behaviour replicated across many
deployments, scaling idle capacity down is a meaningful lever on the energy cost of
everyday computing, which is why the design makes it a default rather than an option.

\paragraph{Societal context.}
The societal concern is the reliability of the services that depend on SmartLoad.
People book medical appointments, sit examinations, and reach essential information
through web applications, and a sustained slowdown during a demand peak is felt as a
real denial of service to real people. This concern places an asymmetric weight on
safety: a system inserted into the request path of such services must never make
them \emph{less} reliable than the classical baseline it replaces. That asymmetry is
the societal origin of the design's single most important property, the
deterministic fallback and operator override on every adaptive decision
(\cref{dsa:sec:reqs-nonfunctional}), and of the decision to keep the decision
plane off the critical request path, so that a failed engine stops adaptation rather
than traffic.

\paragraph{Global context.}
The global concern is portability across heterogeneous deployment environments. An
adaptive middleware that is welded to one cloud provider, one load balancer, or one
orchestration platform cannot serve the diverse global base of teams it is intended
for. The design answers this concern in three concrete commitments that recur as
requirements below: the data plane sits behind a pluggable load-balancer adapter so
that NGINX is a default rather than a lock-in; every external surface speaks a
vendor-neutral standard (\cref{dsa:sec:standards}); and the entire stack ships as a
container composition that an adopter can stand up locally regardless of their host
environment. Portability is thus not an afterthought but a property the standards and
adapter choices are selected to guarantee.

%======================================================================
\section{Requirements and Technical Specifications}
\label{dsa:sec:reqs}

The requirements below are derived directly from the problem statement of
\cref{dsa:sec:problem} and the contextual concerns of \cref{dsa:sec:contexts}. They
are split into functional requirements, meaning what the system must \emph{do},
organised around the capabilities the prototype delivers, and non-functional
requirements, meaning the qualities every part of the system must \emph{exhibit}, of
which safety is first among equals. Each requirement is numbered so that the architecture of
\cref{ch:design_solution} and the evaluation of \cref{ch:discussion} can refer back
to it.

\subsection{Functional Requirements}
\label{dsa:sec:reqs-functional}

\Cref{dsa:tab:fr} enumerates the functional requirements. Each one corresponds to a
capability that is realised as a service or a contract in the architecture, and each
is paired in the rightmost column with the deterministic behaviour it must degrade
to, because, for this system, the fallback is part of the requirement, not a
separate concern.

\begin{table}[ht]
\centering
\caption[Functional requirements and their fallback baselines]{Functional requirements, each naming the deterministic baseline it must degrade to.}
\label{dsa:tab:fr}
\small
\renewcommand{\arraystretch}{1.25}
\begin{tabular}{@{}p{0.06\linewidth}p{0.46\linewidth}p{0.40\linewidth}@{}}
\toprule
\textbf{ID} & \textbf{Requirement} & \textbf{Deterministic baseline / note} \\
\midrule
FR-1 & Ingest per-request telemetry (per-backend latency, error rate, request count)
from the data plane and store it in a time-series database for the decision engines
to query. & Standardised OTLP transport; the request path never blocks on ingestion. \\
FR-2 & Classify each backend as healthy / degraded / unhealthy from its telemetry and
exclude unhealthy backends from the upstream pool, with an audit trail. & Threshold
rule (latency / error-rate cut-offs) is the shipping baseline and the fallback. \\
FR-3 & Produce a short-horizon forecast of incoming request rate and publish it with
confidence bounds. & Moving-average forecaster is the shipping baseline; the trained
forecaster is an opt-in upgrade. \\
FR-4 & Drive autoscaling from the forecast: grow the pool ahead of a predicted peak
and shrink it as demand recedes, respecting a cooldown timer and the policy capacity
bands. & Reactive sliding-window rule takes over when the forecast stream is stale or
absent. \\
FR-5 & Produce learned routing recommendations across the healthy backends, operable
in a passive \emph{shadow} mode (observe only) or an \emph{active} mode (influence
traffic). & Classical scheduler (round-robin / least-connections) is the floor;
shadow is the default and active requires explicit opt-in. \\
FR-6 & Manage a single operating policy (read, validated write, idempotent no-op
detection) and publish every accepted change, recording a per-field audit row. &
Atomic write; a change matching on-disk state produces no audit row and no publish. \\
FR-7 & Present an operator UI exposing system health, the active policy, the audit
trail, manual override actions, and a live view of the decision engines. & The UI is
read-and-control only; it actuates nothing the API does not already expose. \\
FR-8 & Provide an operator override on every adaptive decision: manual backend
isolation, manual scaling, manual routing weights, and a single safe-mode flag that
collapses all adaptation to classical routing. & Overrides bypass the engines but
remain bounded by the policy (e.g.\ manual scale is clamped to the capacity bands). \\
FR-9 & Guarantee a deterministic fallback per engine: if a trained model is disabled,
fails to load, or is overridden, its dependency-free baseline silently takes its
place. & A ``no decision available'' state must never arise; the baseline is always
present. \\
FR-10 & Expose a consolidated status read, one call returning every service's
health, the active policy, and the most recent audit events. & Always returns
\texttt{200}; per-service failures surface inside the body, bounded by a per-service
timeout. \\
\bottomrule
\end{tabular}
\end{table}

Three features of the table warrant emphasis because they encode the central design
commitment. First, FR-2, FR-3, FR-5, and FR-9 each pair an adaptive capability with a
named deterministic baseline; the baseline is not a contingency bolted on afterwards
but a co-requirement specified at the same time as the capability it backs. Second,
FR-5 specifies two operating modes for learned routing precisely because the design
must accommodate operators of differing risk tolerance: shadow mode lets the learned
policy be evaluated against live traffic without influencing it, and the promotion to
active mode is a deliberate operator act, not a default. Third, FR-8 and FR-9 together
express the override hierarchy: an operator can always intervene manually, and even
absent any operator action the system degrades to its baseline rather than to a
failure. These three properties are what make the functional surface defensible in
the request path, and they recur as the safety hierarchy detailed in
\cref{ch:design_solution}.

FR-2 deserves a closer look, because it answers two backend-degradation modes that a
blind scheduler and the load balancer's own health checks both miss.
\Cref{dsa:fig:failclass} sets the two side by side. In the first, a backend carries a
persistent latency penalty, for example a steady extra 15~ms from a
slow disk or a noisy co-tenant, yet keeps answering every request successfully. A
round-robin scheduler assumes every backend is interchangeable, so it keeps handing
the slow instance its full share of traffic, and the slow instance silently drags the
pool's tail latency up for as long as it stays in rotation. In the second, a backend
suffers a sharp latency spike, on the order of an added 200~ms, but
still returns \texttt{200~OK} on every request. NGINX's passive health check trips only
on connection failures and error responses through its \texttt{max\_fails} counter, so
a backend that is merely slow while still succeeding never crosses that counter and is
never excluded. Both mechanisms are invisible to the classical data plane precisely
because nothing fails in the sense the data plane checks for; what changes is latency,
which the scheduler does not measure and the passive check does not watch. FR-2 exists
to close exactly this gap by classifying backends on their latency and error-rate
telemetry and excluding a degraded one that the data plane would otherwise keep
serving.

\begin{figure}[ht]
\centering
\begin{tikzpicture}[
  font=\footnotesize,
  box/.style={slnode,fill=slbox,text width=4.2cm,align=left,inner sep=5pt,minimum height=1.0cm},
  blind/.style={slnode,fill=slhead,text width=4.2cm,align=center,inner sep=4pt},
  miss/.style={slnode,fill=tbdred!18,draw=tbdred,text width=4.2cm,align=center,inner sep=4pt},
  lbl/.style={font=\footnotesize\bfseries},
  arr/.style={slflow,-{Stealth[length=2.5mm]}},
]
  % left column: persistent heterogeneity
  \node[lbl] (la) at (0,0) {(a) Persistent heterogeneity};
  \node[box,below=3pt of la] (a1)
    {Backend carries a steady latency penalty (e.g.\ \texttt{+}15~ms),
     but every request still returns \texttt{200~OK}.};
  \node[blind,below=8pt of a1] (a2)
    {Round-robin assumes interchangeable backends and keeps sending the slow one its
     \(1/N\) share.};
  \node[miss,below=8pt of a2] (a3)
    {Pool tail latency rises and stays high for as long as the backend is in rotation.};
  \draw[arr] (a1) -- (a2);
  \draw[arr] (a2) -- (a3);

  % right column: latency spike, still 200 OK
  \node[lbl] (lb) at (6.4,0) {(b) Latency spike, still \texttt{200~OK}};
  \node[box,below=3pt of lb] (b1)
    {Backend latency spikes (e.g.\ \texttt{+}200~ms) during a window,
     yet responses remain \texttt{200~OK}.};
  \node[blind,below=8pt of b1] (b2)
    {NGINX \texttt{max\_fails} counts only failed responses, so a slow-but-successful
     backend never trips it.};
  \node[miss,below=8pt of b2] (b3)
    {Passive health check never excludes the backend; the spike is served to users.};
  \draw[arr] (b1) -- (b2);
  \draw[arr] (b2) -- (b3);
\end{tikzpicture}
\caption[Two latency-only degradations the data plane misses]{Two backend-degradation mechanisms that round-robin and NGINX passive health checks both miss, since the backend keeps returning \texttt{200~OK} while only its latency degrades.}
\label{dsa:fig:failclass}
\end{figure}

\subsection{Non-Functional Requirements}
\label{dsa:sec:reqs-nonfunctional}

Where the functional requirements say what the system does, the non-functional
requirements in \cref{dsa:tab:nfr} say how it must behave while doing it. Safety
leads the list deliberately: for a request-path component serving the societal role
described in \cref{dsa:sec:contexts}, the quality of \emph{degrading well} dominates
the quality of \emph{performing well}, and the architecture of
\cref{ch:design_solution} is organised around honouring NFR-1 above all others.

\begin{table}[ht]
\centering
\caption[Non-functional requirements, led by safety]{Non-functional requirements, where safety (NFR-1) is the dominant quality.}
\label{dsa:tab:nfr}
\small
\renewcommand{\arraystretch}{1.25}
\begin{tabular}{@{}p{0.07\linewidth}p{0.24\linewidth}p{0.61\linewidth}@{}}
\toprule
\textbf{ID} & \textbf{Quality} & \textbf{Specification} \\
\midrule
NFR-1 & Safety / graceful degradation & Every adaptive decision has a deterministic
fallback and an operator override; loss of any intelligent component reduces the
system to classical routing, never to an outage. The decision plane runs off the
critical request path, so a failed engine stops adaptation rather than traffic. \\
NFR-2 & Low decision latency & Decisions are computed off the request path and
consumed as pre-computed signals; the data plane never queries the database or the
control bus per request, so adaptation adds no synchronous latency to a served
request. \\
NFR-3 & Hot reconfiguration & Routing and capacity changes are applied to the live
data plane without dropping in-flight connections (NGINX upstream rewrite plus
reload), and a reconfiguration must never empty the pool. \\
NFR-4 & Observability & Every service emits structured logs, per-process metrics in a
standard format, and a uniform \texttt{/health} endpoint (a control-plane liveness
signal, not a data-plane serviceability guarantee); every adaptive decision is
attributable and every actuated action is audited. \\
NFR-5 & Testability & Every engine ships with a baseline requiring no model artifact,
decision logic is unit-testable in isolation, and per-request fidelity is verifiable
end to end on live traffic. \\
NFR-6 & Extensibility & Decision engines are pluggable behind a message contract and
load-balancer back-ends are pluggable behind an adapter interface, so a new engine or
a new data plane is a localised change rather than a cross-cutting rewrite. \\
NFR-7 & Portability & The whole stack stands up from a single container composition on
a single host, and every external surface speaks a vendor-neutral standard, so the
design is trivially portable across deployment environments. \\
NFR-8 & Low coupling & The data plane functions with zero intelligent services
running, and no service imports another service's code; services interact only through
the shared contracts. \\
\bottomrule
\end{tabular}
\end{table}

The non-functional requirements are not independent wishes; several reinforce one
another around the safety thesis. NFR-1 (safety) is made achievable by NFR-2 and
NFR-8: because decisions are computed off the request path (NFR-2) and the data plane
is decoupled from the intelligence (NFR-8), a failed engine genuinely cannot take
traffic down; it can only stop the system adapting. NFR-3 (hot reconfiguration)
carries an explicit safety clause, namely that a reconfiguration must never empty the
pool, because dropping the last healthy backend is exactly the kind of edge case a
request-path component must be engineered against from the start; the design's
defence against it is described under the risks of \cref{dsa:sec:risks}. NFR-6
(extensibility) is what lets the design ship a deterministic baseline and a trained
model for the same engine side by side, swapping one for the other behind a stable
contract; it is therefore a precondition for NFR-1 rather than a mere convenience.

%======================================================================
\section{Constraints}
\label{dsa:sec:constraints}

The design operates under real constraints of resource, time, and chosen technology,
and the accreditation context \cite{abet2025criteria} requires these to be named in a
recognised vocabulary rather than left implicit. \Cref{dsa:tab:constraints} maps the
project's constraints onto that vocabulary, distinguishing the constraints imposed
\emph{on} the project (cost, schedule) from those the design \emph{adopts as
commitments} (standards, maintainability, extensibility, interoperability, usability,
safety) and those \emph{imposed by the chosen tools}.

\begin{table}[ht]
\centering
\caption[Project constraints by accreditation category]{Project constraints, mapped onto the accreditation constraint vocabulary \cite{abet2025criteria}.}
\label{dsa:tab:constraints}
\small
\begin{tabular}{@{}p{0.22\linewidth}p{0.72\linewidth}@{}}
\toprule
\textbf{Constraint} & \textbf{How it bears on the design} \\
\midrule
Cost / economic & No paid components: the entire stack is built from open-source data,
control, and observability layers, so an adopter incurs no licence cost and the team
incurs none building it. This rules out, for example, the paid NGINX module that would
otherwise provide dynamic upstream weights, forcing the configuration-rewrite-and-reload
actuation path the design adopts instead. \\
Schedule & One academic year, a team of three. This bounds the achievable surface and
is the direct reason multi-tenancy, authentication, and the managed control plane are
deferred to Phase~2 (\cref{dsa:sec:problem-scope}) rather than attempted thinly. \\
Standards & The external surfaces are constrained to the established standards of
\cref{dsa:sec:standards}; the design may not invent private wire formats where a
vendor-neutral standard exists. \\
Maintainability & Cross-service interactions must be typed, documented contracts,
namely a single shared module for message envelopes and a single shared module for SQL,
so that the system remains comprehensible and changeable by a small team; ad-hoc JSON and
scattered query strings are prohibited by convention and checked in CI. \\
Extensibility & Decision engines and load-balancer back-ends must be replaceable behind
their contracts (NFR-6), so that a research engine under active iteration can be swapped
without touching the data plane or the autoscaler. \\
Interoperability & The system must compose with, not replace, the surrounding
observability and data-plane ecosystem: it consumes a standard load balancer and emits
standard telemetry, so it slots into an existing estate rather than demanding one. \\
Usability & A single operator must be able to inspect health, change policy, read the
audit trail, take manual actions, and watch the engines live from one UI, without
reading source or shelling into containers. \\
Safety & The overriding constraint: every adaptive decision is bounded by a
deterministic fallback and an operator override (NFR-1). No design choice may silently
remove that floor. \\
Tooling-imposed & The chosen tools impose their own constraints: a single-node Docker
Compose deployment means no genuine multi-host scaling in the prototype; an in-memory
Redis bus means message loss must be assumed and tolerated; and actuation through an
NGINX reload means reconfiguration granularity is bounded by reload cadence rather than
being instantaneous. \\
\bottomrule
\end{tabular}
\end{table}

Two constraints deserve a closing word because they are felt most sharply in the
chapters that follow. The schedule constraint, one year and three people, is the
reason the project draws a clear line between the adaptive control loop (built and
demonstrated) and the SaaS surface (deferred), and it is a stated constraint that we
own openly: the deferred items are named and planned, not silently dropped. The
tooling-imposed constraints are the reason several of the risks in
\cref{dsa:sec:risks} are worth naming at all. The single-node deployment bounds the
scale at which results can be measured, and the in-memory bus is the reason every
consumer of an adaptive signal is required, by NFR-1, to tolerate that signal's
absence.

%======================================================================
\section{Risks and Mitigations}
\label{dsa:sec:risks}

A strong design chapter records the risks it anticipates and shows, for each, the
mitigation already built into the design rather than left to chance. The point is not
to list dangers but to demonstrate that the architecture was shaped around them from
the start. \Cref{dsa:tab:risks} lists the principal risks with their likelihood,
impact, and mitigation; the prose that follows expands on the two that bear most
directly on the request path, since the design's handling of them is a direct
demonstration of the safety thesis in action.

\begin{table}[ht]
\centering
\caption[Principal project risks and their mitigations]{Principal project risks, each paired with a mitigation built into the design.}
\label{dsa:tab:risks}
\small
\renewcommand{\arraystretch}{1.25}
\begin{tabular}{@{}p{0.27\linewidth}p{0.09\linewidth}p{0.09\linewidth}p{0.43\linewidth}@{}}
\toprule
\textbf{Risk} & \textbf{Likelihood} & \textbf{Impact} & \textbf{Mitigation} \\
\midrule
A learned policy might not improve on a strong deterministic baseline on every
workload. & Medium & Medium &
The deterministic fallback (FR-9) guarantees the established behaviour is always
available, and the learned router ships in shadow mode (FR-5) behind that baseline; it
is promoted to active only on demonstrated benefit, so routing is always at least as
good as the proven approach while the policy is refined. \\
A model could over-react and exclude healthy capacity. & Medium & High &
The deterministic threshold rule is the default for guaranteed stability (FR-2), with
the trained model offered as an enhanced opt-in; a quorum guard in the data plane
refuses any exclusion that would drop the last active backend (NFR-3); and non-finite
telemetry features are filtered before scoring. \\
Telemetry pipeline gaps could corrupt downstream decisions. & Medium & High &
Standardised OTLP transport (\cref{dsa:sec:standards}); per-request fidelity is asserted
end to end by an integration test on live traffic; the request path never blocks on
ingestion (NFR-2). \\
A single-node Compose deployment bounds the scale at which the system can be evaluated.
& High & Low & Accepted as a prototype constraint (\cref{dsa:sec:constraints}); the
design stays trivially portable to a multi-host orchestrator (NFR-7), and the contracts
already assume message loss and horizontal replication. \\
Evaluation evidence could lag implementation while the learned components are measured.
& Medium & Medium & The deterministic baselines let the system ship safely while the
learned components are characterised, so this is a bounded measurement task on named
components, not a structural unknown. \\
A forecast could be wrong and point the autoscaler in the wrong direction. & Medium & Medium &
Confidence bounds accompany every forecast (FR-3); the reactive sliding-window fallback
(FR-4) bounds the effect of a bad prediction; a cooldown timer prevents oscillation. \\
\bottomrule
\end{tabular}
\end{table}

\paragraph{Routing risk: a learned policy may not beat a strong baseline everywhere.}
The most consequential routing risk is that a learned policy will not improve on a
strong deterministic baseline on every workload; on a heterogeneous mix, plain
round-robin is already a tough opponent. The design answers this risk by construction
rather than by hope. Because FR-9 guarantees a deterministic fallback and FR-5 pins
the learned router in shadow mode by default, the learned policy runs behind a proven
classical scheduler and a fallback at all times. The classical scheduler carries
traffic; the learned policy observes and records what it would have done; and an
operator promotes it to active only once it demonstrates a clear benefit on the target
workload. Routing is therefore always at least as good as the established approach
while the learned policy is refined, which is exactly the guarantee the safety
thesis sets out to provide. The full quantitative comparison of the learned router
against the classical baselines is reported in \cref{ch:discussion}; the modelling
choices behind the policy are discussed in \cref{ch:litreview}.

\paragraph{Exclusion risk: keeping an over-reacting engine from emptying the pool.}
The second routing risk is sharper because its impact is high: an anomaly engine that
over-reacts under sustained load could classify a healthy backend as unhealthy, and in
the limit exclude the last one in the pool. This is exactly the case NFR-3's ``must
never empty the pool'' clause is written to cover, and the design answers it on three
layers. First, the conservative threshold rule is the shipping default, chosen for
guaranteed stability so the stack is safe out of the box, with the trained engine
available as an explicit, enhanced opt-in. Second, and structurally, a quorum guard in
the data plane refuses any exclusion that would drop the last active backend, removing
the empty-pool case at its root and protecting the threshold engine as well as the
trained one. Third, non-finite telemetry features are filtered before scoring, so
spurious inputs cannot drive an unhealthy verdict. A safe, stable default plus a
structural guard that no engine can bypass is the design's safety discipline applied
to the most dangerous failure mode a request-path component can have, and it is the
clearest demonstration in this chapter that the deterministic-fallback commitment is
real rather than rhetorical.

\paragraph{Anticipated risks and the evidence schedule.}
The remaining rows of \cref{dsa:tab:risks} are forward-looking, and one deserves a word
of its own: evaluation evidence may lag implementation. The control loop is
demonstrated end to end and the deterministic baselines are solid, while the
head-to-head quantitative case for several learned components is still being assembled.
The design's answer to this risk is the same property that answers the others: because
every learned component degrades to a measured, dependable baseline, the evidence
schedule is a bounded measurement task on named components rather than an open question
about whether the system works at all. The full quantitative treatment of these
results, including the comparison of the learned router against the classical
baselines, is reported in \cref{ch:discussion}; here it is enough to record that the
design ships safely while that evidence is gathered, and that this, too, is the safety
thesis doing its work.
